\section{Discusión y hallazgos propios del grupo}
\label{sec:discusion}

\subsection{Contradicciones y vacíos}

Las fuentes revisadas no presentan contradicciones directas sobre la necesidad de migrar hacia criptografía poscuántica, pero sí muestran diferencias importantes entre las recomendaciones estratégicas y el grado de madurez práctica de la transición. Joseph et al. recomiendan adoptar \textit{crypto-agility} mediante abstracciones, bibliotecas criptográficas centralizadas y sistemas preparados para cambiar algoritmos con bajo impacto \cite{joseph2022transitioning}. Sin embargo, la revisión sistemática de Näther et al. encuentra que la literatura todavía no posee una definición uniforme de \textit{crypto-agility} y que las propuestas concretas para implementarla son heterogéneas \cite{naether2024migrating}. Por tanto, un vacío relevante no es reconocer su importancia, sino convertirla en prácticas y criterios suficientemente claros para evaluar si un sistema es realmente \textit{crypto-agile}.

Los mecanismos híbridos ilustran otra tensión. Joseph et al. los presentan como una estrategia para reducir el riesgo de depender inmediatamente de algoritmos poscuánticos relativamente nuevos, combinando componentes clásicos y poscuánticos \cite{joseph2022transitioning}. Näther et al. confirman que estas soluciones aparecen con frecuencia en las migraciones estudiadas, pero también recogen el riesgo de ataques de \textit{downgrade}, en los que un adversario intenta desviar la comunicación hacia algoritmos clásicos \cite{naether2024migrating}. Esto no invalida el enfoque híbrido, pero muestra que la compatibilidad hacia atrás debe diseñarse con cuidado: mantener una alternativa clásica puede facilitar la transición y, al mismo tiempo, conservar una ruta vulnerable si la negociación o selección de algoritmos no está protegida adecuadamente.

% Aporte pendiente del Integrante 5:
% Añadir aquí, si corresponde, contradicciones o vacíos identificados en CISA/NSA/NIST
% y Westerbaan sin repetir la descripción ya incluida en el estado del arte.

\subsection{Conexión con el curso}

Los hallazgos se relacionan directamente con la tríada de
\textbf{confidencialidad, integridad y disponibilidad (CID)}. En
confidencialidad, el riesgo \textit{store now, decrypt later} descrito por
Joseph et al. muestra que información capturada hoy podría ser descifrada
en el futuro ante avances en computación cuántica
\cite{joseph2022transitioning}. FIPS 203 responde a esta amenaza mediante
ML-KEM, que permite establecer secretos compartidos para su posterior uso
con criptografía simétrica \cite{nist2024fips203}.

En integridad y autenticidad, FIPS 204 y FIPS 205 estandarizan ML-DSA y
SLH-DSA como mecanismos de firma digital poscuántica
\cite{nist2024fips204,nist2024fips205}. Joseph et al. destacan además la
importancia de proteger las firmas en sistemas difíciles de actualizar,
pues estas pueden utilizarse posteriormente para distribuir actualizaciones
confiables \cite{joseph2022transitioning}.

La disponibilidad también puede verse afectada durante la transición.
Näther et al. identifican mayores requisitos de hardware y rendimiento,
además de la necesidad de mantener interoperabilidad entre sistemas
\cite{naether2024migrating}. Por ello, una migración que introduzca
incompatibilidades o degradaciones importantes puede mejorar la resistencia
criptográfica y, simultáneamente, afectar la disponibilidad.

Finalmente, los roles de migración propuestos por Näther et al., como
\textit{Migration Manager}, \textit{Asset \& Risk Manager},
\textit{Security Expert} y \textit{Developer/Administrator}, pueden
relacionarse con los temas de control de acceso y asignación de
responsabilidades vistos en el curso \cite{naether2024migrating}. Aunque
no constituyen un modelo RBAC formal, muestran una separación de
responsabilidades dentro del proceso de migración.

% Aporte pendiente del Integrante 5:
% Reforzar la conexión con CID, controles o políticas usando sus fuentes.

\subsection{Interpretación del grupo}

Consideramos que disponer de algoritmos poscuánticos estandarizados no
implica que una organización esté preparada para la transición. El reto
se desplaza hacia la capacidad de identificar dónde se utiliza
criptografía vulnerable, priorizar los sistemas según su riesgo y
dependencias, y sustituir mecanismos sin comprometer la interoperabilidad
o disponibilidad. Por ello, interpretamos que la preparación para PQC
debe evaluarse no solo por los algoritmos adoptados, sino también por la
capacidad de la infraestructura para adaptarse a futuros cambios
criptográficos \cite{joseph2022transitioning,naether2024migrating,  
nist2024fips203,nist2024fips204,nist2024fips205}.

% Aporte pendiente del Integrante 5:
% Integrar aquí los hallazgos derivados de CISA/NSA/NIST y Westerbaan.
